% SEGMENT OLP-0512-S01
% Part: intuitionistic-logic
% Chapter: tableaux
% Section: introduction

\documentclass[../../../include/open-logic-section]{subfiles}

\begin{document}

\olfileid{int}{tab}{int}

\olsection{அறிமுகம்}

!!^{tableau}s என்பவை கீழ்நோக்கிக் கிளைப்படையும் சில மரங்கள்; அவை
!!{signed formula}s, அதாவது மெய்மதிப்புக் குறி ($\True$ அல்லது $\False$)
மற்றும் !!a{sentence} ஆகியவற்றால் ஆன சோடிகளைக் கொண்டவை:
\[
\sFmla{\True}{!A} \text{ அல்லது } \sFmla{\False}{!A}.
\]
!!^a{tableau} பல \emph{எடுகோள்களுடன்} தொடங்குகிறது. அதற்குப் பின்வரும்
ஒவ்வொரு !!{signed formula} உம் உய்த்தறி விதிகளுள் ஒன்றைப் பயன்படுத்தி
உருவாக்கப்படுகிறது. சில உய்த்தறி விதிகள் மரத்தின் ஒரு நுனியில் ஒன்று அல்லது
அதற்கு மேற்பட்ட !!{signed formula}s ஐச் சேர்க்கின்றன; மற்றவை இரண்டு புதிய
நுனிகளைச் சேர்த்து இரு கிளைகளை உருவாக்குகின்றன. விதி பயன்படுத்தப்பட்ட
!!{signed formula} இன் !!{formula} ஐவிடச் சிக்கல் குறைந்த வாய்பாடுகள்
கொண்ட !!{signed formula}s ஐ விதிகள் உருவாக்குகின்றன. ஒரு கிளை
$\sFmla{\True}{!A}$, $\sFmla{\False}{!A}$ ஆகிய இரண்டையும் கொண்டிருந்தால்,
அக்கிளை \emph{மூடியது} என்கிறோம். !!a{tableau} இன் ஒவ்வொரு கிளையும்
மூடியிருந்தால், முழு !!{tableau} உம் மூடியது. மூடிய !!{tableau} ஒன்று,
!!{tableau} ஐத் தொடங்கப் பயன்படுத்திய !!{signed formula}s இன் கணம்
நிறைவுறத்தக்கதல்ல என்று காட்டும் !!a{derivation} ஆகும். இதைப் பயன்படுத்தி
$\Proves$ உறவை வரையறுக்கலாம்: $\Gamma \Proves !A$ எனில், எனில் மட்டுமே,
$\Gamma_0 = \{!B_1, \dots, !B_n\} \subseteq \Gamma$ ஆகும் ஏதோவொரு
முடிவுறு கணத்திற்குப் பின்வரும் எடுகோள்களுக்கான மூடிய !!{tableau} ஒன்று
உள்ளது:
\[
\{\sFmla{\False}{!A}, \sFmla{\True}{!B_1}, \dots, \sFmla{\True}{!B_n}\}.
\]

உள்ளுணர்வுவாதத் தருக்கத்திற்காக, !!{signed formula} என்ற கருத்தை
விரிவாக்குவதுடன் இணைப்பிகளுக்கான விதிகளையும் மாற்றியமைக்க வேண்டும். ஒரு
குறியுடன் ($\True$ அல்லது $\False$), இவ்வகை !!{tableau}s இலுள்ள
!!{formula}s க்கு \emph{முன்னொட்டுகள்}~$\sigma$ உம் உள்ளன.
% SOURCE CORRECTION TA-IL-017: this chapter concerns intuitionistic, not modal, tableaux.
முன்னொட்டுகள் நேர்ம முழு எண்களின் காலியல்லாத தொடர்கள்; அதாவது
$\sigma \in (\PosInt)^* \setminus \{\emptyseq\}$. அத்தகைய முன்னொட்டுகளைச்
சூழ்ந்துள்ள $\tuple{\ }$ இல்லாமல் எழுதி, தனித்தனி !!{element}s ஐ
$,$ களுக்குப் பதிலாக~$.$ களால் பிரிக்கிறோம். $\sigma$ ஒரு முன்னொட்டு
எனில், $\sigma.n$ என்பது $\sigma \concat \tuple{n}$; எடுத்துக்காட்டாக,
$\sigma = 1.2.1$ எனில் $\sigma.3$ என்பது $1.2.1.3$. ஆகவே,
\[
\sFmla{\True}{!A \lif (!B \lif !C)}[1.2]
\]
என்பது ஒரு \emph{முன்னொட்டுள்ள !!{signed formula}} (சுருக்கமாக ஒரு
\emph{முன்னொட்டுள்ள !!{formula}}).

உள்ளுணர்வுப்படி, முன்னொட்டு என்பது !!a{tableau} கிளையிலுள்ள !!{formula}s ஐ
நிறைவுசெய்யக்கூடிய மாதிரி ஒன்றின் உலகிற்குப் பெயரிடுகிறது; $\sigma$ ஓர்
உலகிற்குப் பெயரிடுகிறது எனில், $\sigma.n$ ஆனது~$\sigma$ பெயரிடும்
உலகிலிருந்து அணுகத்தக்க ஓர் உலகிற்குப் பெயரிடுகிறது.

உள்ளுணர்வுவாத மாதிரிகளில் அணுகுறவு தற்சுட்டுப் பண்பும் கடப்புப் பண்பும்
கொண்டது. முன்னொட்டுகளின் வகையில், $\sigma$ ஆனது~$\sigma$ இலிருந்தே
அணுகத்தக்கது;
மேலும் $\sigma$ வை நீட்டிக்கும் எந்த முன்னொட்டும், அதாவது
$\sigma.n_1.\cdots.n_k$ வடிவிலுள்ள எந்த முன்னொட்டும், அணுகத்தக்கது என்பது
இதன் பொருள். $\sigma$ வையும் அதன் எந்த நீட்டிப்பையும் குறிக்க $\sigma.*$
என்ற குறிமுறையை அறிமுகப்படுத்துவோம். வேறு சொற்களில், $\sigma.*$
முன்னொட்டுகள்,~$\sigma$ இலிருந்து அணுகத்தக்க முன்னொட்டுகள் அனைத்தும், அவை
மட்டுமே.

\end{document}
